% Classical supplementary proof for G2-8; not an assembled Paper III.
% GPT-6 Astra (OpenAI, ChatGPT Pro mode), effort Pro, 17 November 2026.
\begin{proposition}[The orientation-compatible almost complex homotopy class]
\label{mapsaudit:almost-complex-homotopy}
Any two almost complex structures on the tangent bundle of the smooth
six-sphere that induce the same orientation are homotopic through almost
complex structures. Consequently, after an orientation-preserving
diffeomorphism $d:S^6\to X$, the pulled-back structure $d^*J_X$ is
homotopic to the octonionic almost complex structure with that orientation.
The resulting complex rank-three tangent bundles are isomorphic.
\end{proposition}
\begin{proof}
Fix a Riemannian metric $g$. For any almost complex structure $J$, define
the positive definite form
\[
 h_J(v,w)=\tfrac12\bigl(g(v,w)+g(Jv,Jw)\bigr)=g(H_Jv,w).
\]
The positive self-adjoint endomorphism $H_J$ depends smoothly on $J$,
and $h_J$ is $J$-invariant. Hence
\[
 J_t=H_J^{t/2}JH_J^{-t/2}\quad(0\leq t\leq1)
\]
is a smooth path of almost complex structures with $J_1$ orthogonal for
$g$. If $J$ is already orthogonal, $H_J=I$ and the path is constant.
Every conjugating endomorphism has positive determinant, so the
orientation is preserved. It suffices to compare orthogonal structures.

On an oriented Euclidean six-space their space is
$SO(6)/U(3)$: an orthonormal complex basis identifies any structure
with the standard one, and its stabilizer is exactly $U(3)$.
Here is a direct identification with $\mathbb{CP}^3$ that also specifies
which component is being used. Give $\mathbb C^4$ its standard Hermitian
form and unit volume form. The associated antilinear Hodge star on
$\bigwedge^2\mathbb C^4$ has square one. Its fixed vectors form a real
Euclidean six-space $W$, and
$W\otimes\mathbb C=\bigwedge^2\mathbb C^4$.
The exterior-square action gives a homomorphism
\[
                    SU(4)\longrightarrow SO(W).
\]
It is onto. Indeed, diagonalizing an element of the Lie algebra as
$\operatorname{diag}(i\theta_1,\ldots,i\theta_4)$ shows that it acts
trivially on the exterior square only if every
$\theta_j+\theta_k=0$, and thus every $\theta_j=0$.
The derivative is therefore injective between Lie algebras of equal
dimension fifteen. The image is an open and closed subgroup of the
connected group $SO(6)$, hence all of it. The same eigenvalue argument
shows that the group kernel is $\{I,-I\}$.

For a line $L\subset\mathbb C^4$ put
$U_L=L\wedge L^\perp$. The star interchanges $U_L$ and
$\bigwedge^2L^\perp$, so each vector of $W$ has a unique expression
$u+*u$ with $u\in U_L$. Define
\[
                  J_L(u+*u)=iu-i*u.
\]
This is an orthogonal complex structure. The construction is
$SU(4)$-equivariant and gives a smooth map from $\mathbb{CP}^3$ onto
one orientation component of the orthogonal structures: transitivity
of $SU(4)$ on lines and surjectivity onto $SO(6)$ prove surjectivity.
It is injective because $U_L$, its $i$-eigenspace, recovers $L$ as
\[
 L=\{v\in\mathbb C^4: v\wedge u=0\text{ for every }u\in U_L\}.
\]
For example, when $L=\mathbb Ce_1$, testing $e_1\wedge e_j$ for
$j=2,3,4$ proves this identity coefficient by coefficient.
Equivalently the two homogeneous spaces have the same stabilizer
$S(U(1)\times U(3))$ in $SU(4)$, so this bijection is a
diffeomorphism. Choosing the orientation of $W$ induced by one
$J_L$ identifies the desired component; the other is obtained by
reversing the orientation.

The Hopf fibration $S^1\to S^7\to\mathbb{CP}^3$ gives
\[
             \pi_6(\mathbb{CP}^3)=\pi_6(S^7)=0.
\]
Here the two adjacent homotopy groups of $S^1$ vanish, and
$\pi_6(S^7)=0$ follows from cellular approximation. The fibre
$\mathbb{CP}^3$ is also path connected.

The bundle of orientation-compatible orthogonal structures on $TS^6$
is an associated bundle with this fibre. It is not being assumed to
be a product. Use the CW decomposition of $S^6$ with one zero-cell
and one six-cell, and pull the bundle back to $S^6\times[0,1]$.
The given sections prescribe its values at the two ends. Choose a
path between their values over the zero-cell. Over the characteristic
six-disc the pulled-back bundle is trivial. The prescribed data on
\[
 (D^6\times\{0,1\})\cup(\partial D^6\times[0,1])
                    =\partial(D^6\times[0,1])\cong S^6
\]
give a map into the fibre. It extends over $D^6\times[0,1]$ because
$\pi_6(\mathbb{CP}^3)=0$. Compatibility on the characteristic
boundary is exactly the prescribed path over the zero-cell, so the
extension descends and is the required section homotopy. There are
no intermediate cells and no further obstruction groups.
Smooth approximation in bundle charts, followed by a tubular
retraction to the smooth fibre and performed relative to the two
ends, makes this homotopy smooth.

A homotopy of complex structures gives an isomorphism of the end
complex bundles. Explicitly, for a smooth path $J_t$ solve the
linear bundle ODE
\[
 \dot A_t=\tfrac12J_t\dot J_t A_t,\qquad A_0=I.
\]
The identity $J_t\dot J_t+\dot J_tJ_t=0$ gives
$[\tfrac12J_t\dot J_t,J_t]=\dot J_t$; uniqueness of the ODE
implies $A_tJ_0A_t^{-1}=J_t$. Thus $A_1$ is the asserted complex
bundle isomorphism. Apply this to $d^*J_X$ and the octonionic
structure with its orientation. If the labelled octonionic field
in the working course induces the opposite orientation, it is its
negative that enters this assertion.
\end{proof}

\paragraph{Credit and scope.}
This is a classical consequence of the twistor fibre and elementary
section obstruction theory. It was stated in the reader's note G2-8,
with an abbreviated argument. The proof above fills the homogeneous-space
identification and the nontrivial-bundle extension step; it does not
claim any new integrability theorem or a fibration-adapted
diffeomorphism $d$.
